% ch5_c §5.8–5.9 (书页 166–170 / p180–p184)
%--B-TAIL--
%JOIN%场所测得的那样——我们可以写出介质的一个\emph{复折射率}（complex refractive index）$\mathbf{N}$：
\begin{equation*}
\mathbf{N}^{2}=\epsilon(\mathbf{q},\omega)\approx 1-\frac{\omega_{p}^{2}}{\omega^{2}}. \tag{5.67}
\end{equation*}
若 $\omega<\omega_{p}$，则 $\epsilon$ 为负，$\mathbf{N}$ 为纯虚数，这引起光的\emph{全反射}。但若 $\omega>\omega_{p}$，则 $\epsilon$ 为正，$\mathbf{N}$ 为实数；当沿垂直于表面的方向观察时，金属就变成\emph{透明}的。当对电子引入散射机制之后，这个理论要稍作修改（见 §8.6），但在碱金属中确实观察到了\emph{紫外透明}现象。

\section{准粒子与内聚能}

当然，上述基于介电常数的分析只是一种近似。实质性的步骤在 §5.1 中已经完成：在那里假定势的每一个 Fourier 分量都可以独立地处理。这在使用术语上称为\emph{无规相位近似}（Random Phase Approximation），在对这个问题几乎所有更细致的分析中都不得不采用这一假定。

除了作一些一般性的评论之外，我们将不讨论这些更复杂的理论。结果发现，与 Coulomb 势相联系的电子之间的强长程相互作用，可以约化成类似式 (5.29) 的屏蔽相互作用。电子不仅在杂质周围形成电荷云、把远距离处的场屏蔽掉；每个电子还可以说都随身携带着自己的电荷云——实际上是一个“正”的云，对应于把“其他”电子有效地排除在它的邻域之外。

反过来，在晶体中到处运动的\emph{正电子}（positron）会把一个“负”的电荷云吸引到自己周围。因此，金属中\emph{正电子湮灭}（positron annihilation）的速率，由于正电子中心处电子概率密度的增强，要比未受扰动的均匀电子气情形略高一些。但是，多体效应并不会严重干扰对湮灭过程中所发射的两个 $\gamma$ 射线量子方向之间角关联的简单解释。由于认为正电子在湮灭之前已在晶格中静止下来，发射方向对相反方向的任何偏离，必定来自湮灭电子所提供的动量（图 96）。因此，角 $\theta$ 的分布给出了传导电子气中动量分布的一种量度，这或许会被证明是对金属与合金电子结构理论的一种有用检验。

%FIG{96}: {正电子湮灭中的动量守恒}

严格说来，一个带着正电荷晕的电子，需要用非常复杂的波函数才能得到对它的恰当描述。然而，它的行为确实非常像一个粒子，因为它携带着电子电荷 $e$。“其他”电子的效应，是改变这个量子态的能量与波矢之间的关系；例如，可以有
\begin{equation*}
\mathcal{E}(\mathbf{k})=\frac{\hbar^{2}k^{2}}{2m^{*}}, \tag{5.68}
\end{equation*}
其中 $m^{*}$ 并不完全等于普通的自由电子质量 $m$。结果发现，这些准粒子（quasi-particle）在费米能级附近是相当稳定的，但在远离费米面时便趋于瓦解和耗散。这一点对于输运理论以及费米面概念的形式论证十分重要。

准粒子图像实际上正是电子气的\emph{元激发}（elementary excitations）谱的一种表示。虽然我们不能对真实金属精确地算出这个谱，但对于处在正凝胶（§5.1）中的、具有相应密度的理想电子气，我们可以对它的性质作出相当好的估计。例如，这个模型的基态能量就告诉我们一些关于金属内聚能（cohesive energy）（§4.3）的信息。Wigner 和 Seitz 的粗糙假定——电子把它到访的原胞中的所有其他传导电子统统驱赶出去——可以换成对电子与其屏蔽云之间静电相互作用的直接计算，从而对金属键合中电子关联（electron correlation）的效应作出好得多的估计。

同样，我们应当预期两个准粒子之间的相互作用相当弱，多少就如同通过屏蔽的 Coulomb 相互作用那样。因此，我们的系统符合\emph{Fermi 液体}（Fermi liquid）的 Landau 模型：其中较低的诸能级用把一个或多个准粒子激发到更高的态来表示，就好像我们拥有一团极低温度下的由\emph{独立}费米子（fermions）组成的气体一样。但是，当我们把越来越多的粒子激发到费米能级之上时，就必须计及它们之间的相互作用。这些相互作用由下列形式的附加能量项表示：
\begin{equation*}
E_{2}=\frac{1}{2}\sum_{\mathbf{k},\sigma;\,\mathbf{k}',\sigma'}f(\mathbf{k},\sigma;\,\mathbf{k}',\sigma')\,\delta n(\mathbf{k},\sigma)\,\delta n(\mathbf{k}',\sigma') \tag{5.69}
\end{equation*}
其中 $\delta n(\mathbf{k},\sigma)$ 是当我们从基态过渡到系统的某个较高能级时，波矢为 $\mathbf{k}$、自旋为 $\sigma$ 的准粒子模式的占据数的变化。

Landau 模型的参数决定了传导电子气的谱，从而决定了诸如电子比热（§4.7）这样的一些性质。相互作用项 $f(\mathbf{k},\sigma;\mathbf{k}',\sigma')$ 的一个重要特点是：它确实依赖于两个电子的自旋的相对取向。在普通金属密度下，这种\emph{交换}（exchange）力有利于平行排列——这本质上是因为 Pauli 原理使两个自旋相同的电子不致占据空间同一点，从而增强了关联效应并稍稍降低了它们的相互静电势能。这个效应并不大，但显然在金属磁性的理论中起着重要作用（见 §§10.2--10.5）。

\section{Mott 转变}

还有一种本质上起因于电子之间作用力的定性效应，它并不常被观察到，但在原理上颇为重要。设想我们把大量的氢原子（或 Na 原子——或者任何每个原胞含有奇数个电子的原子阵列）缓慢地聚拢在一起。按照 §4.1 中阐述的原理，每个原子上的电子波函数将发生交叠，能带将会形成，固体就应当立即变成导体。而且看起来，无论原子相距多远，这都应当发生。

这是一个悖论性的结果，难以令人相信。实际上，很可能并非如此。§3.4 中曾提示，一组局域轨道（localized orbitals）$\phi_{a}(\mathbf{r}-\mathbf{l})$ 与一组 Bloch 函数 $\phi_{\mathbf{k}}(\mathbf{r})$ 形式上等价；但这种等价性只有在我们忽略 Pauli 原理和电子间相互作用时才成立。如果每个原子有两个电子，那么满带中所有态 $\phi_{\mathbf{k}}(\mathbf{r})$ 构成的行列式波函数，与所有局域原子轨道\footnote{严格说来，这里我们应当指的是 Wannier 函数（见 §6.2）。但是，假如我们假定当原子相距很远时紧束缚近似成立，论证也完全一样。}$\phi_{a}(\mathbf{r}-\mathbf{l})$ 构成的行列式相同。但如果每个原子只有一个电子，半满带中诸态的行列式就不等同于局域轨道的行列式。它倾向于把两个电子安放到某些原子上，而让另一些原子近乎空着——这样的电荷分布可能给系统的总 Coulomb 能贡献一些很大的正项。因此，作为系统的试探波函数，局域态的集合也许比离域（delocalized）的能带态集合更好。前一类波函数描述绝缘体，后一类描述导体。

为了区分这两者，我们可以采用如下的论证。设想从局域态阵列中取出一个电子。它将在身后留下一个空的离子，即一个正电荷。电子会被这个离子吸引，甚至可能在它那里形成一个束缚态，以致它实际上并不能自由地穿越晶体去传导电流。但假如已经有许多这样的电子存在，它们就会屏蔽离子的电荷，此时离子的势可以取为如下形式：
\begin{equation*}
-\frac{e^{2}}{r}\,e^{-\lambda r}, \tag{5.70}
\end{equation*}
如同式 (5.29) 一样。若已电离电子的密度为 $n$，则按照式 (5.22)，
\begin{equation*}
\lambda^{2}\approx\frac{4me^{2}n^{1/3}}{\hbar^{2}} \tag{5.71}
\end{equation*}
（利用了式 (3.3)、(3.4) 和 (4.3)）。

普通的 Coulomb 场是能够束缚一个电子的。最低的态——氢原子的基态——具有半径
\begin{equation*}
a_{H}=\frac{\hbar^{2}}{me^{2}}. \tag{5.72}
\end{equation*}
如果屏蔽半径 $1/\lambda$ 小于 $a_{H}$，那么连这个态也无法被束缚。因此，具有势 (5.70) 的电离原子能够重新俘获从它那里取走的电子，其条件是
\begin{gather*}
\lambda<a_{H}^{-1}, \\
\text{即}\quad n^{-1/3}>4a_{H}. \tag{5.73}
\end{gather*}
如果我们原子的平均间距大于 4 个原子单位，金属态就不能自我维持，系统应当变成绝缘体。而且，这是一种合作效应；每一个额外的自由电子都有助于使其余的电子松脱。即使我们没有把式 (5.73) 中的数字 4 算对，这个转变在原则上看来也会相当陡峭；把我们的氢或 Na 原子聚拢起来，我们应当会看到它们在某个确定的原子间距处突然从绝缘体转变为金属。

\emph{Mott 转变}（Mott transition）已经在某些过渡金属氧化物中观察到，它们在某个温度以上非常突然地变成导电的。这个理论也可以应用于半导体中的杂质中心，此时须对 $a_{H}$ 加以修正，即乘以半导体的体介电常数（如 §5.6 中那样；原书误印作 §5.5）（见 §6.4）。

上述论证其实与 \emph{Wigner 假说}密切相关：密度极低的电子气将倾向于“结晶”。令 $r_{0}$ 为一个电子所占据的球的平均半径，即
\begin{equation*}
\tfrac{4}{3}\pi r_{0}^{3}=1/n. \tag{5.74}
\end{equation*}
每个电子的势能，由于在它周围有一圈带正电的背景凝胶球，将是
\begin{equation*}
-e^{2}/2r_{0}. \tag{5.75}
\end{equation*}
的量级。但每个电子由于被限制在这个球的区域之内，还会有量级为
\begin{equation*}
\hbar^{2}/mr_{0}^{2} \tag{5.76}
\end{equation*}
的零点能（zero-point energy）。因此，当
\begin{gather*}
e^{2}/2r_{0}\geqslant\hbar^{2}/mr_{0}^{2}, \\
\text{即}\quad r_{0}\geqslant 2a_{H}. \tag{5.77}
\end{gather*}
时，势能将占上风。其中的数值因子，如同式 (5.73) 中的一样，不必当真；而且在什么样的条件下这个转变实际上可能被观察到，也完全不清楚。尽管如此，认识到第 4 章的简单能带理论以及本章的简单电子关联理论的这些局限性，是十分重要的。
